% Generated draft: requires mathematical and source review.
\documentclass[11pt]{article}
\usepackage[margin=25mm]{geometry}
\usepackage{fontspec}
\setmainfont{Latin Modern Roman}
\usepackage{amsmath,amssymb,mathtools,unicode-math}
\setmathfont{Latin Modern Math}
\usepackage{enumitem,needspace,xurl,hyperref}
\setlength{\parindent}{0pt}
\setlength{\parskip}{5pt}
\setlength{\emergencystretch}{3em}
\newcommand{\N}{\mathbb{N}}
\newcommand{\Z}{\mathbb{Z}}
\newcommand{\Q}{\mathbb{Q}}
\newcommand{\R}{\mathbb{R}}
\newcommand{\C}{\mathbb{C}}
\newcommand{\F}{\mathbb{F}}
\newcommand{\A}{\mathbb{A}}
\newcommand{\PP}{\mathbb{P}}
\newcommand{\E}{\mathbb{E}}
\newcommand{\eps}{\varepsilon}
\newcommand{\dd}{\,\mathrm{d}}
\begin{document}
\section*{UnsolvedMath ID: 30005711}
\textbf{Source problem code:} catalogue entry
\subsection{Definitions and mathematical statement}
Let \(d,n\in\N\) satisfy \(d\geq 3\) and \(n\geq d+1\). Let \(K_n\) be the complete simple graph on \(n\) vertices, and put \(m=\binom{d+1}{2}\). For a real symmetric matrix \(A\), write \(\lambda_i(A)\) for its \(i\)-th smallest eigenvalue, counting multiplicity.

For a graph \(G=(V,E)\) and a placement \(p:V\to\R^d\), the supplied excerpt uses a rigidity matrix \(R(G,p)\) with \(d|V|\) rows and defines its stiffness matrix by
\[
L(G,p)=R(G,p)R(G,p)^{\mathsf T}\in\R^{d|V|\times d|V|}.
\]
The excerpt states the rank bound
\[
\operatorname{rank}R(G,p)\leq d|V|-\binom{d+1}{2}
\]
and calls a framework infinitesimally rigid when equality holds. It calls a graph \(d\)-rigid when some placement yields an infinitesimally rigid framework, and asserts equivalence with infinitesimal rigidity at generic placements. In the present range \(|V|=n\geq d+1\), the displayed bound makes \(\lambda_{m+1}(L(G,p))\) the eigenvalue immediately following the guaranteed \(m\) zero eigenvalues.

The source attributes to Jord\'an and Tanigawa the definition
\[
a_d(G)=\sup_{p:V\to\R^d}\lambda_{m+1}\bigl(L(G,p)\bigr).
\]
\textbf{Unresolved definition.} The supplied excerpt does not give the entries or normalization of \(R(G,p)\). Its displayed supremum ranges over all maps, but it does not explain whether coincident vertex positions are admissible or how the rigidity matrix is defined there. The separate observation about dimension one assumes an injective placement and identifies the rigidity matrix with an incidence matrix; this does not supply the missing higher-dimensional definition. Thus the formula above retains the source's notation and quantifier without supplying an unsupported normalization or placement restriction.

\textbf{Requested conjecture, conditional on recovering those conventions.} With \(a_d\) interpreted exactly as intended in the source, establish
\[
\forall d\in\N\;\forall n\in\N,\qquad
\bigl(d\geq3\ \text{and}\ n\geq d+1\bigr)
\ \Longrightarrow\ a_d(K_n)\leq\frac nd.
\]
This is the explicit target in the supplied curated statement. The original excerpt instead says that a displayed lower bound is conjectured to be essentially tight; that qualitative assertion alone does not specify the exact upper bound above. Moreover, the extracted lower bound contains the corrupted expression \(a_d(K_n)\geq\text{``21 nd''}\), so its coefficient and fraction structure cannot be recovered reliably from this record alone.
\subsection{Short English statement}
For every integer dimension \(d\geq3\) and every \(n\geq d+1\), the supplied curated problem asks whether the \(d\)-dimensional algebraic connectivity of \(K_n\) is at most \(n/d\). A complete formulation requires confirmation of the rigidity matrix's normalization and its admissible placements. The supplied original excerpt expresses only conjectured essential sharpness of a lower bound.
\subsection{Sources}
\begin{itemize}
\item Local UnsolvedMath registry, permanent ID \(30005711\), title \emph{Algebraic Connectivity of Complete Graphs}, source problem code \texttt{catalogue entry}. The supplied record labels the statement \texttt{corrected\_verified}, says it was reconstructed and checked against the official report, and labels the problem \texttt{partially\_solved}. These are supplied curation claims, not independent verification here.
\item \emph{Geometric, Algebraic and Topological Combinatorics} (2024), Oberwolfach Reports, pp.\ 3249--3318, DOI \texttt{10.4171/owr/2023/58}: \url{https://doi.org/10.4171/owr/2023/58}. The extracted passage identifies the contribution \emph{Rigidity expander graphs}, Eran Nevo, joint work with Alan Lew, Yoval Peled, and Orit Raz, and includes the page marker 3283. The registry records proposed year 2023.
\item The excerpt attributes the definition of higher-dimensional algebraic connectivity to Jord\'an and Tanigawa, citing its internal reference [5]. The bibliographic entry for [5] is not supplied. Its internal reference [1], cited for rigidity facts, is also not identified in the record.
\item Lew--Nevo--Peled--Raz, \emph{Rigidity Expander Graphs}: \url{https://doi.org/10.1007/s00493-025-00149-z}.
\item Lew--Nevo--Peled--Raz, \emph{algebraic connectivity of graphs}: \url{https://arxiv.org/abs/2205.05530}.
\item Supplied literature assessment, checked on \textbf{2026-08-22}: the displayed lower bound is reported proved, its claimed essential sharpness is reported conjectural, and later work is reported not to establish matching asymptotics. This is supplied dated curation derived from the OWR record, not an independently verified literature assessment; the record requests independent verification of time-sensitive status and cited claims.
\end{itemize}
\end{document}
